\chapter{Grouped and latent-group models}
\label{ch:templates}
\chapterpromise{Known-group shared-changepoint analysis and a finite latent-group criterion with Jensen minorization and exact template updates over a finite partition set.}

\section{Known groups}
Suppose each sequence $s$ has an observed group label $g(s)\in\{1,\ldots,G\}$. Within a group,
sequences may share a common partition while retaining sequence-specific segment parameters. Under
conditional independence, the group-specific marginal likelihood for a candidate segment is the
product of the sequence-specific segment marginal likelihoods. The sum-product and max-sum
recursions of Chapters~\ref{ch:sumproduct} and~\ref{ch:map} then apply separately to each group.
This is a direct hierarchical extension of the shared-changepoint model in Chapter~\ref{ch:shared}.

\section{Finite latent-group criterion}
When group labels are unknown, let $\tau_g=(k_g,t^{(g)})$ denote a deterministic candidate
changepoint template for group $g$. The archived implementation uses the sequence--template score
\begin{equation}
 S_s(\tau)=\frac{p(k)}{C_k}\prod_{q=1}^{k}
 A^{(0,s)}_{t_{q-1},t_q}\,\gamma_u(t_{q-1},t_q),
\label{eq:template-sequence-score}
\end{equation}
with the convention that an impossible segment gives score zero. This subject-template score is not assumed
to integrate to one over the sample space of sequence $s$. For weights $\pi_g\ge0$ with
$\sum_g\pi_g=1$, define
\begin{equation}
 \mathcal F(\pi,\tau)=\sum_{s=1}^{S}\log\left\{\sum_{g=1}^{G}\pi_g S_s(\tau_g)\right\}.
\label{eq:template-lik}
\end{equation}
Equation~\eqref{eq:template-lik} is the optimization criterion studied in this chapter. The results
below concern this finite criterion and do not require interpreting $S_s(\tau_g)$ as a normalized
mixture component.

\begin{figure}[H]
\centering
\resizebox{0.96\textwidth}{!}{\begin{tikzpicture}[node distance=8mm and 9mm]
\node[secondarybox,text width=31mm] (scores) {Sequence-specific template scores\\$S_s(\tau_g)>0$};
\node[primarybox,right=of scores,text width=29mm] (resp) {Auxiliary allocation weights\\$r_{sg}$};
\node[primarybox,right=of resp,text width=32mm] (update) {Alternating updates\\$\pi_g$ and $\tau_g$};
\node[successbox,right=of update,text width=31mm] (obj) {$\mathcal F=\sum_s\log\sum_g\pi_gS_s(\tau_g)$};
\draw[arrPrimary] (scores)--(resp);
\draw[arrPrimary] (resp)--(update);
\draw[arrPrimary] (update)--(obj);
\draw[arrSecondary] (obj.north) |- ++(0,7mm) -| node[pos=0.52,above,font=\scriptsize]{next iteration} (scores.north);
\node[neutralbox,below=10mm of resp,text width=101mm] {The updates maximize a Jensen minorizer of the stated finite latent-group criterion. The criterion is not interpreted as a normalized finite-mixture likelihood.};
\end{tikzpicture}
}
\caption{Alternating optimization of the finite latent-group criterion. Auxiliary allocation
weights make a Jensen lower bound tight; group-template updates maximize the resulting weighted
partition score over the finite candidate set.}
\label{fig:latent-group}
\end{figure}

\section{Jensen minorization and alternating updates}
\label{sec:latent-em}
For any nonnegative weights $r_{sg}$ satisfying $\sum_g r_{sg}=1$, the log-sum inequality gives
\begin{equation}
\mathcal F(\pi,\tau)\ge
\mathcal Q(r,\pi,\tau):=
\sum_{s,g}r_{sg}\left[\log\pi_g+\log S_s(\tau_g)-\log r_{sg}\right].
\label{eq:template-bound}
\end{equation}
Equality holds at
\begin{equation}
 r_{sg}=\frac{\pi_gS_s(\tau_g)}{\sum_h\pi_hS_s(\tau_h)}.
\label{eq:template-resp}
\end{equation}
For fixed $r$, maximization over the group weights gives
$\pi_g=S^{-1}\sum_s r_{sg}$. Let $n_g=\sum_s r_{sg}$. Substituting
Equation~\eqref{eq:template-sequence-score} into $\mathcal Q$ shows that, for a candidate segment
count $k$, the group-$g$ template contribution is
\begin{equation}
 n_g\{\log p(k)-\log C_k\}
 +\sum_{q=1}^{k}\ell^{(g)}_{t_{q-1},t_q},
\qquad
 \ell^{(g)}_{ij}=\sum_{s=1}^{S}r_{sg}\log A^{(0,s)}_{ij}
                  +n_g\log\gamma_u(i,j).
\label{eq:template-blockscore}
\end{equation}
Thus the update is an exact max-sum recursion at fixed $k$, followed by comparison across $k$ using
the displayed count offset. The convention $0\log0=0$ applies when $r_{sg}=0$, while a positive
responsibility on a zero sequence--template score makes that candidate value $-\infty$.

\begin{theorem}[Monotonicity of exact latent-group updates]
\label{thm:em-monotone}
Consider an iteration that (i) sets $r$ according to Equation~\eqref{eq:template-resp}, (ii)
maximizes $\mathcal Q$ over $\pi$, and (iii) exactly maximizes each responsibility-weighted template
criterion using the local score and segment-count offset in Equation~\eqref{eq:template-blockscore}
over its declared finite candidate set. Then the iteration cannot decrease $\mathcal F$.
\end{theorem}
\begin{proof}
At the current parameter values, Equation~\eqref{eq:template-resp} makes the Jensen bound tight:
$\mathcal F_{\rm old}=\mathcal Q(r,\pi_{\rm old},\tau_{\rm old})$. Exact maximization of the
minorizer gives
\[
\mathcal Q(r,\pi_{\rm new},\tau_{\rm new})\ge
\mathcal Q(r,\pi_{\rm old},\tau_{\rm old}).
\]
The log-sum inequality gives
$\mathcal F_{\rm new}\ge\mathcal Q(r,\pi_{\rm new},\tau_{\rm new})$.
Combining the inequalities yields $\mathcal F_{\rm new}\ge\mathcal F_{\rm old}$.
\end{proof}

\begin{algorithm}[H]
\caption{Alternating optimization for finite latent-group templates}
\label{alg:multi-em}
\KwIn{Sequence-specific segment-score arrays; number of groups $G$; admissible segment counts; number of restarts; convergence tolerance.}
\ForEach{restart}{
 initialize $\pi$ and templates $\tau_{1:G}$\;
 evaluate the current value of $\mathcal F$\;
 \Repeat{the increase in $\mathcal F$ is at most the tolerance or the iteration limit is reached}{
  update $r_{sg}$ by Equation~\eqref{eq:template-resp}\;
  update $\pi_g\leftarrow S^{-1}\sum_s r_{sg}$\;
  for each group $g$, set $n_g\leftarrow\sum_s r_{sg}$, build the local scores in Equation~\eqref{eq:template-blockscore}, solve the max-sum recursion at every allowed $k$, and select the offset-adjusted maximizer using $n_g\{\log p(k)-\log C_k\}$\;
  evaluate and append the new value of $\mathcal F$\;
 }
 retain the restart with the largest final value in its recorded objective sequence\;
}
\KwOut{Group templates, group weights, auxiliary allocation weights, complete objective sequences, and restart ranking by the final objective value.}
\end{algorithm}

\section{Non-uniqueness and implementation limitation}
The criterion is invariant to permutation of the group labels. It may have duplicate or collapsed
templates when $G$ is larger than the effective number of groups, and multiple restarts do not
guarantee a global maximum. These are optimization and identifiability limitations of the stated
criterion.

The archived implementation contains an exit path that can return the objective value from the
previous iteration rather than the last value in the objective sequence. In one of three diagnostic
restarts, the difference was $0.003137360860080$. The archived latent-group simulation remains a
real computation, but \taskid{CODE-BB-003} must correct the returned objective and restart ranking
before those values are treated as release-quality optimization results.

\begin{statusbox}[Established result and non-claim]
For Equation~\eqref{eq:template-lik}, the auxiliary allocation update makes the Jensen minorizer
tight, the group-weight update maximizes that minorizer, and the finite template step is solved
exactly by max-sum recursion when Equation~\eqref{eq:template-blockscore} applies. The monotonicity
theorem is complete. No theorem is claimed for identifiability of a normalized finite mixture or
consistency of the estimated number of groups.
\end{statusbox}
